\documentclass[12pt,a4paper]{article}

% Pacotes essenciais
\usepackage[utf8]{inputenc}
\usepackage[brazil]{babel}
\usepackage[T1]{fontenc}
\usepackage{times}
\usepackage{graphicx}
\usepackage{indentfirst}
\usepackage{setspace}
\usepackage{geometry}
\usepackage{titlesec}
\usepackage{tocloft}
\usepackage[hidelinks]{hyperref}
\usepackage{amsmath}
\usepackage{amssymb}

% Configuração das margens (ABNT)
\geometry{
    left=3cm,
    right=2cm,
    top=3cm,
    bottom=2cm
}

% Espaçamento entre linhas
\onehalfspacing
\sloppy
\emergencystretch=5em

% Configuração de seções
\titleformat{\section}{\normalfont\normalsize\bfseries\MakeUppercase}{\thesection}{1em}{}
\titleformat{\subsection}{\normalfont\normalsize\bfseries}{\thesubsection}{1em}{}
\titleformat{\subsubsection}{\normalfont\normalsize}{\thesubsubsection}{1em}{}

% Configuração do sumário
\renewcommand{\cftsecfont}{\normalfont}
\renewcommand{\cftsecpagefont}{\normalfont}
\renewcommand{\cftsecleader}{\cftdotfill{\cftdotsep}}

% Comandos personalizados
\newcommand{\nomealuno}{PETERSON CARARA JUNIOR}
\newcommand{\titulo}{Formalização de Máquinas de Estados Finitos Componíveis em Haskell}
\newcommand{\subtitulo}{Segurança de Tipos e Completude Funcional via Monads e Tipos Algébricos}
\newcommand{\curso}{Curso Superior de Tecnologia em Ciência de Dados}
\newcommand{\orientador}{Prof. Felipe Canarozzo Lourenço}
\newcommand{\cidade}{Santos - SP}
\newcommand{\ano}{2025}

\begin{document}

% ========== CAPA ==========
\begin{titlepage}
    \begin{center}
        \MakeUppercase{\textbf{Centro Estadual de Educação Tecnológica Paula Souza}}\\
        \MakeUppercase{\textbf{Faculdade de Tecnologia da Baixada Santista ``Rubens Lara''}}\\
        \MakeUppercase{\textbf{\curso}}\\
        \MakeUppercase{\textbf{Programa de Iniciação Científica}}
        
        \vspace{6cm}
        
        \MakeUppercase{\textbf{\nomealuno}}
        
        \vspace{4cm}
        
        \MakeUppercase{\textbf{\titulo: \subtitulo}}
        
        \vfill
        
        \cidade\\
        \ano
    \end{center}
\end{titlepage}

% ========== FOLHA DE ROSTO ==========
\begin{titlepage}
    \begin{center}
        \MakeUppercase{\textbf{\nomealuno}}
        
        \vspace{6cm}
        
        \MakeUppercase{\textbf{\titulo: \subtitulo}}
        
        \vspace{3cm}
        
        \begin{flushright}
            \begin{minipage}{8cm}
                \begin{spacing}{1.0}
                    Projeto de pesquisa apresentado à Faculdade de Tecnologia da Baixada Santista ``Rubens Lara'', do \curso, no Programa de Iniciação Científica.
                    
                    \vspace{0.5cm}
                    
                    Orientador: \orientador
                \end{spacing}
            \end{minipage}
        \end{flushright}
        
        \vfill
        
        \cidade\\
        \ano
    \end{center}
\end{titlepage}

% ========== LISTA DE ABREVIATURAS ==========
\newpage
\section*{LISTA DE ABREVIATURAS E SIGLAS}
\addcontentsline{toc}{section}{LISTA DE ABREVIATURAS E SIGLAS}

\begin{tabular}{ll}
ABNT & Associação Brasileira de Normas Técnicas\\
ADT & Algebraic Data Type (Tipo de Dados Algébrico)\\
CREM & Crème de la Crem: Composable Representable Executable Machines\\
DDD & Domain-Driven Design\\
DFSM & Deterministic Finite State Machine\\
FATEC BS & Faculdade de Tecnologia da Baixada Santista Rubens Lara\\
FSM & Finite State Machine (Máquina de Estados Finitos)\\
I/O & Input/Output\\
PFP & Programação Funcional Pura\\
\end{tabular}

% ========== SUMÁRIO ==========
\newpage
\tableofcontents

% ========== CORPO DO TEXTO ==========
\newpage
\section{INTRODUÇÃO}

As Máquinas de Estados Finitos (FSMs) constituem uma das abstrações mais duradouras e essenciais da Ciência da Computação. Elas fornecem o formalismo matemático rigoroso para a modelagem de sistemas cujo comportamento é estritamente determinado por um conjunto discreto de estados e regras de transição. O modelo é crucial para a especificação do comportamento de sistemas reativos, abrangendo desde protocolos de rede até a lógica de domínio em arquiteturas de software complexas. Conforme Perone e Karachalias (2023), é possível construir arquiteturas de software a partir da composição de máquinas de estados.

No entanto, a tradução desse modelo abstrato para linguagens de programação de propósito geral do paradigma imperativo, como C++, é frequentemente marcada pela dependência de variáveis mutáveis, tipicamente campos de classe ou variáveis globais, para armazenar o estado atual. Essa prática compromete a pureza do modelo matemático. Segundo Malakhovski e Shalyto (2009), a implementação em linguagens imperativas não permite a aplicação direta dos métodos tradicionais de FSMs em um ambiente funcional. A dependência da mutabilidade e de tipos fracos para o estado expõe o sistema a erros lógicos e transições inválidas, delegando a responsabilidade de verificar a corretude à fase de testes em tempo de execução.

O presente projeto de pesquisa propõe resolver essa discrepância fundamental através da formalização e implementação de FSMs em Haskell. Por ser uma linguagem puramente funcional (Hutton, 2016), Haskell exige que a função de transição seja implementada como uma função pura. O diferencial reside no seu sistema de tipos avançado: a utilização de Tipos de Dados Algébricos (ADTs) é essencial para a modelagem de dados (Milewski, 2020), permitindo que os conjuntos finitos de estados e eventos sejam mapeados diretamente para tipos de dados. Essa correspondência de tipo não é apenas uma convenção, mas uma imposição que, por meio da verificação de exaustividade do \textit{pattern matching}, viabiliza a verificação da correção das funções de transição no estágio de compilação (Malakhovski \& Shalyto, 2009).

Além da segurança de tipos básica, a abordagem em Haskell permite a exploração de conceitos avançados para o tratamento do estado. A capacidade de sequenciar transições e lidar com o estado de forma pura é resolvida através de Monads. Essa abstração, que expressa de forma concisa programas que manipulam estado (Wadler, 1992), permite o tratamento puro do estado através de estruturas como a \texttt{State Monad} (Milewski, 2020). Esse arcabouço monádico, estruturado pela Categoria Kleisli (Maia, 2022), fornece o formalismo necessário para a composição segura de transições.

A capacidade de composabilidade é o ponto culminante desta abordagem. O \textit{framework} Crem (Crème de la Crem: Composable Representable Executable Machines), desenvolvido por Perone e Karachalias (2023), ilustra como os \textit{advanced type-level features} de Haskell podem ser usados para especificar transições de estado permitidas e proibidas. Essa especificação transforma as regras de domínio em propriedades intrínsecas ao tipo de dados, permitindo a composição de máquinas de estados independentes (Perone e Karachalias, 2023), resultando em um sistema robusto e modular, onde a lógica de domínio é validada pelo compilador antes da execução.

\subsection{PROBLEMA DE PESQUISA}

Considerando o descompasso entre a formalidade matemática das FSMs e a fragilidade de tipos nas implementações tradicionais (Malakhovski \& Shalyto, 2009), o problema de pesquisa é formulado como:

\begin{quote}
   De que forma a linguagem de programação funcional pura Haskell, explorando seu sistema de tipos avançados (Tipos de Dados Algébricos) e suas abstrações de ordem superior (Monads e Categoria Kleisli), pode ser utilizada para obter uma formalização de Máquinas de Estados Finitos que ofereça garantias verificáveis de correção (\textit{type safety}) e completude (\textit{completeness}) na especificação das funções de transição, e que suporte a composição de  FSMs complexas  de maneira modular e representável? 
\end{quote}

\subsection{JUSTIFICATIVA DO TEMA}

A relevância do projeto reside em transferir a verificação da corretude da FSM da fase de testes para a fase de compilação. Segundo Malakhovski e Shalyto (2009), uma das vantagens da implementação funcional é a verificação da correção do código das funções de transição já no estágio de compilação, o que é possível pelo uso de Tipos de Dados Algébricos (ADTs) para representar eventos e estados.

Além disso, a estrutura de Haskell é baseada na Teoria das Categorias (Milewski, 2020). Essa base formal permite lidar com o estado de forma pura através de Monads, que foram originalmente desenvolvidas por teóricos de categorias para expressar concisamente certos aspectos da álgebra universal (Wadler, 1992). Formalmente, uma monad é definida como um Monóide na Categoria dos Endofuntores (Jahn, 2020) e permite a mudança de estado de forma pura (Milewski, 2020). Essas estruturas possibilitam a composição segura de transições via a Categoria Kleisli (Maia, 2022; Jahn, 2020). Essa composabilidade é crucial, pois permite a construção de arquiteturas de software a partir da composição de máquinas de estados independentes (Perone e Karachalias, 2023), o que contribui para a modularidade e a coerência do domínio do \textit{software}.

\subsubsection{Hipóteses}
O sistema de tipos de Haskell oferece uma Segurança de Tipos superior, onde o uso dos Tipos de Dados Algébricos (ADTs) e o \textit{pattern matching} conseguem impor, em tempo de compilação, a proibição de estados e transições logicamente impossíveis. Essa garantia se estende à Completude Funcional, onde o \textit{pattern matching} exaustivo permite verificar que a função de transição cubra todas as combinações de (Estado $\times$ Evento), o que, segundo Malakhovski e Shalyto (2009), seria uma vantagem sobre as linguagens imperativas; essa vantagem será avaliada na Fase 4. Finalmente, postula-se a Hipótese de Composabilidade e Modulação: a aplicação de abstrações categoriais como Monads e recursos avançados de \textit{type-level} de Haskell (como os utilizados na biblioteca Crem) permite a criação de Máquinas de Estados Finitos Componíveis de forma puramente funcional, resultando em uma arquitetura modular cujas regras de domínio são validadas pelo compilador (Perone e Karachalias, 2023).

\subsubsection{Objetivo Geral}
O objetivo geral desta Iniciação Científica é formalizar e implementar Máquinas de Estados Finitos Determinísticas (DFSMs) em Haskell, explorando o sistema de tipagem avançado e as estruturas funcionais da linguagem para garantir a segurança e a completude na definição das funções de transição, com ênfase na demonstração de um modelo que suporte a composição de sistemas.

\subsubsection{Objetivos Específicos}

Para alcançar o objetivo geral de formalizar e implementar FSMs de forma segura e componível em Haskell, o projeto se desdobra nos seguintes objetivos específicos.

O primeiro objetivo visa o alicerce do projeto: definir modelos de dados em Haskell que traduzam com precisão os elementos formais de uma FSM. Esta formalização será realizada através da criação de Tipos de Dados Algébricos (ADTs), nos quais os conjuntos finitos de Estados ($Q$) e Eventos ($\Sigma$) serão mapeados diretamente para Tipos Soma. Essa técnica visa incorporar o rigor do modelo matemático ao código, seguindo o princípio de que ADTs são o conceito mais importante na modelagem funcional (Milewski, 2020).

O segundo objetivo trata da validação da lógica interna da máquina: implementar o mecanismo de processamento de sequências de entrada e a função de transição ($\delta$). Serão exploradas as funções de alta ordem e, sobretudo, o \textit{pattern matching} exaustivo para definir a função $\delta: \text{Estado} \times \text{Evento} \to \text{Estado}$. Essa abordagem é fundamental para atingir a Hipótese de Completude Funcional, pois o compilador de Haskell verifica a exaustividade (como aviso, que o projeto promove a erro de compilação), de modo que todo par $(\text{Estado}, \text{Evento})$ seja tratado, transferindo a verificação da correção do código das funções de transição para o estágio de compilação (Malakhovski \& Shalyto, 2009).

Em seguida, o terceiro objetivo foca na gestão do estado e na composabilidade: explorar a aplicação do \textit{State Monad} e da Categoria Kleisli. Será demonstrado como a \texttt{State Monad} permite gerenciar a sequência de mudanças de estado de forma pura, encapsulando o estado e garantindo a mudança de estado de forma pura (Milewski, 2020). Esta etapa é crítica para lidar com o histórico de transições. Além disso, a Categoria Kleisli será utilizada como base teórica para compor transições monádicas, fornecendo o formalismo para encadear logicamente as transições de estado, como sugerido por Milewski (2020) e Maia (2022).

Finalmente, o quarto objetivo é a validação prática da abordagem: demonstrar a composabilidade de máquinas de estados por meio de um estudo de caso. Este estudo de caso ilustrará como a tipagem de Haskell pode ser usada para garantir transições válidas e o encapsulamento do domínio, permitindo que a lógica do \textit{software} seja estruturada como uma composição de máquinas de estados modulares (Perone e Karachalias, 2023). O resultado esperado é um sistema onde as regras de transição --- inclusive a especificação de transições de estado permitidas e proibidas (Perone e Karachalias, 2023) --- são validadas pelo compilador, confirmando o potencial das estruturas funcionais para o desenvolvimento de \textit{software} robusto.

\section{REFERENCIAL TEÓRICO INICIAL}
O presente referencial teórico estabelece a base de conhecimento necessária para a formalização das Máquinas de Estados Finitos (FSMs) no ambiente funcional puro. A abordagem é construída sobre a fusão de quatro domínios conceituais: a Teoria de Autômatos, a Programação Funcional Pura (PFP), a Teoria das Categorias e as Arquiteturas de \textit{Software} Componíveis.

\subsection{Máquinas de Estados Finitos e a Questão da Mutabilidade}

As FSMs são o modelo canónico para descrever sistemas sequenciais e reativos, cuja dinâmica é regida por uma quíntupla formal $(Q, \Sigma, \delta, q_0, F)$, onde $Q$ é o conjunto finito de estados, $\Sigma$ é o alfabeto de entrada (eventos), $\delta$ é a rigorosa função de transição, $q_0$ é o estado inicial, e $F$ é o conjunto de estados finais. A função $\delta: Q \times \Sigma \to Q$ é o coração do sistema, definindo como o autômato se move entre os estados.

A implementação tradicional desse modelo em linguagens de programação imperativas, como C++ ou Java, cria uma divergência perigosa entre o modelo e o código. A representação do estado recorre, invariavelmente, a variáveis mutáveis --- tipicamente um campo de classe (se a programação orientada a objetos for usada) ou uma variável global (Malakhovski \& Shalyto, 2009). Além disso, os identificadores de estado são frequentemente representados por abstrações de baixo nível, como \textit{strings} ou inteiros, que carecem de garantias de tipo. Essa dependência de estado mutável e a fragilidade de tipos resultam numa tradução que não consegue aplicar diretamente os métodos utilizados na implementação de máquinas de estados finitos em linguagens de programação imperativas no rigoroso domínio funcional (Malakhovski \& Shalyto, 2009), expondo a FSM a estados logicamente impossíveis ou transições inválidas que só são detectadas em tempo de execução, através de pesados testes unitários.

\subsection{Haskell e a Garantia de Corretude em Tempo de Compilação}
O projeto propõe a Programação Funcional Pura (PFP) como a solução para este descompasso. Haskell, uma linguagem puramente funcional, baseia-se na imutabilidade e na ausência de efeitos colaterais. Sua filosofia exige que a função de transição ($\delta$) seja implementada como uma função pura, uma transformação que recebe o estado e o evento e retorna, deterministicamente, o novo estado, preservando a imutabilidade do sistema e produzindo \textit{software} claro, conciso e correto (Hutton, 2016).

\subsubsection{Tipos de Dados Algébricos (ADTs)}
A segurança formal em Haskell é fornecida por seu sistema de tipos avançado, ancorado nos Tipos de Dados Algébricos (ADTs). Os ADTs permitem que os conjuntos finitos $Q$ e $\Sigma$ da FSM sejam mapeados diretamente para Tipos Soma, que são considerados o conceito mais importante para a modelagem em Haskell (Milewski, 2020). Por exemplo, se uma FSM possui estados \textit{Ligado} e \textit{Desligado}, o Tipo Soma data Estado = Ligado | Desligado traduz essa finitude para o sistema de tipos. Essa codificação evita o uso de strings e inteiros como identificadores de estado (Malakhovski \& Shalyto, 2009), garantindo que o compilador imponha a integridade estrutural da máquina.

\subsubsection{Pattern Matching e Completude Funcional}
A função de transição é escrita por casamento de padrões (\textit{pattern matching}): cada equação diz qual é o resultado para uma forma possível dos argumentos (Hutton, 2016). Como estado e evento são tipos soma com número finito de construtores, o compilador consegue listar todos os pares (Estado $\times$ Evento) e conferir se cada um tem uma equação. O GHC faz essa conferência, mas o resultado é um aviso, não um erro: quando falta um par, emite o aviso \texttt{-Wincomplete-patterns}, que já vem ligado em \texttt{-Wall} e informa qual par ficou sem tratamento. O aviso só impede a compilação se o projeto pedir, com \texttt{-Werror} ou, de forma mais precisa, \texttt{-Werror=incomplete-patterns}. Com GHC 9.6.7, sem a opção o programa compila e o aviso aparece; com a opção, a compilação falha no mesmo ponto.

A completude funcional depende, portanto, de duas condições: a opção de erro deve estar ligada, e a função não pode usar uma equação curinga (\texttt{\_}), que cobre qualquer par restante e esconde as omissões. Malakhovski e Shalyto (2009, p. 21) afirmam que, com tipos algébricos, o compilador verifica a completude e a ausência de contradições na função de transição; no exemplo deles, uma equação supérflua faz o compilador emitir um aviso.

\subsection{Teoria das Categorias: Monads, Estado Puro e a Categoria Kleisli}
Em um ambiente funcional puro, onde as transições de estado devem ser implementadas como funções puras, o desafio inerente é a gestão da sequência de transições ao longo do tempo. O projeto aborda este desafio recorrendo à Teoria das Categorias, que é a disciplina matemática que fornece os fundamentos lógicos para a linguagem Haskell. O princípio orientador é que a composição é a essência da programação (Milewski, 2020), e a Teoria das Categorias define o arcabouço para a composição segura e coerente.

\subsubsection{Monads: Origem Categorial, Leis e a Abstração do Efeito de Estado}

As Monads são o dispositivo arquitetónico que permite a Programação Funcional Pura gerir efeitos colaterais, como o estado, exceções ou I/O, sem violar a imutabilidade. Segundo Wadler (1992), as mônadas permitem combinar a flexibilidade de recursos impuros com as garantias da pureza, ao encapsular o efeito.

A sua origem é estritamente matemática: Monads foram inventadas na década de 1960 por teóricos de categorias para expressar concisamente certos aspectos da álgebra universal (Wadler, 1992). A definição categorial formal e mais abstrata estabelece uma Monad como um Monóide na Categoria dos Endofuntores (Jahn, 2020). Toda mônada deve obedecer a três leis, identidade à esquerda, identidade à direita e associatividade (Hutton, 2016, cap. 12). Bocchino (2025, seções 4.4 a 4.6) mostra que elas equivalem a exigir que a composição de Kleisli seja associativa e tenha identidade. O compilador não verifica essas leis; elas precisam ser demonstradas para cada instância, o que é feito para a State Monad na seção 5.7.

Para a FSM, a \texttt{State Monad} é a Monad específica empregada. Ela resolve o dilema do estado ao modelar a transição como uma função de transformação de estado, na qual a função de transição pura $\delta$ é encapsulada. Isto permite que a mudança de estado de forma pura seja realizada (Milewski, 2020), pois a função não \textit{modifica} o estado, mas sim \textit{calcula} o novo estado a partir do antigo, retornando o par (Valor, Novo Estado) e mantendo a imutabilidade fora da abstração monádica.

\subsubsection{A Categoria Kleisli para a Composição de Transições}

Embora Monads resolvam o problema do estado encapsulado, a composição dessas funções monádicas (as transições de estado) não pode ser realizada pela composição de funções padrão. É aqui que entra a Categoria Kleisli.

A Kleisli é uma construção categorial específica que é projetada para compor morfismos monádicos, ou seja, funções "sujas" que carregam o efeito do \textit{Monad}. Em vez de funções $A \to B$, a Kleisli compõe morfismos do tipo $A \to MB$ (Maia, 2022), onde $M$ é o Functor Monádico (a \texttt{State Monad} no nosso caso) que carrega o estado. Essa categoria fornece as regras e, mais importante, o operador de composição monádica (\texttt{>=>}) que automatiza a logística da sequenciação.

O uso da Categoria Kleisli é crucial para o projeto, pois ele fornece o formalismo para encadear múltiplas transições de estado da FSM de forma segura e elegante. A Monad, sob a estrutura da Kleisli, garante que o estado do autômato seja automaticamente passado, transformado e extraído em cada etapa da sequência de eventos. Isso simplifica a sintaxe e, desde que a instância da mônada cumpra as três leis, faz com que a composição de transições monádicas seja associativa e tenha identidade (Bocchino, 2025, seção 7), pré-requisito para compor FSMs complexas em arquiteturas modulares.

\subsection{Composabilidade Arquitetural e o Framework CREM}

A validação final da abordagem reside na sua capacidade de suportar arquiteturas de \textit{software} complexas. O CREM nasce de ideias do \textit{Domain-Driven Design} (DDD), em que partes do domínio de um sistema, como agregados, políticas e projeções, são modeladas como máquinas de estados (Perone e Karachalias, 2023, seção 2). Neste trabalho, o interesse recai sobre a composição dessas máquinas, e não sobre o DDD em si.

O \textit{framework} Crème de la Crem (CREM) é um exemplo prático que demonstra o potencial máximo da PFP nesse contexto. O CREM utiliza os recursos avançados de \textit{type-level} de Haskell para permitir que o \textit{software} seja estruturado como uma composição de máquinas de estados independentes (Perone e Karachalias, 2023). Essa modularidade permite que subsistemas de domínio sejam desenvolvidos de forma isolada e testada, e combinados de forma segura.

A maior contribuição do CREM é a capacidade de codificar regras de negócio como tipos. Ele permite ao usuário especificar transições de estado permitidas e proibidas diretamente na assinatura de tipos (Perone e Karachalias, 2023). Isso significa que as restrições mais complexas da lógica de domínio da FSM são transformadas em restrições de tipo, e uma transição que viole a topologia resulta em erro de compilação, e não em erro de \textit{runtime}; regras que dependem de sequências de eventos ou do evento que dispara a transição não são cobertas (seção 5.12). Além disso, o CREM é representável no sentido dos autores: a implementação de um sistema se mantém sincronizada com a sua representação gráfica, pois o diagrama é gerado a partir do próprio código (Perone e Karachalias, 2023). O termo não se refere ao functor representável da teoria das categorias.

\section{PROCEDIMENTOS METODOLÓGICOS}

A metodologia deste projeto será de natureza mista, combinando o rigor da pesquisa formal com a validação prática, pautando-se numa abordagem qualitativa para analisar a fundo o fenômeno da formalização de Máquinas de Estados Finitos (FSMs) no contexto da programação funcional pura. O trabalho se desenvolverá em duas vertentes principais: a Pesquisa Bibliográfica e a Pesquisa Aplicada (ou Tecnológica). A Pesquisa Bibliográfica fornecerá a base formal, coletando e analisando o referencial teórico sobre FSMs, Programação Funcional Pura (Haskell), Tipos de Dados Algébricos (ADTs), Teoria das Categorias (Monads, Categoria Kleisli) e os princípios do \textit{framework} CREM. Esta base teórica é crucial para a construção do modelo computacional.

A Pesquisa Aplicada será a fase de experimentação, focada no desenvolvimento de um artefato de \textit{software} (código em Haskell) para comprovar as hipóteses de segurança e composabilidade. Os procedimentos serão organizados em quatro fases interligadas que visam, em última instância, cumprir todos os objetivos específicos do projeto.

A Fase 1: Revisão Teórica e Formalização Inicial visa estabelecer o alicerce do código. Inicialmente, será concluída a revisão aprofundada dos fundamentos da FSM e de como a imutabilidade da PFP se traduz numa alternativa superior ao modelo imperativo. Em seguida, o primeiro objetivo específico será atendido pela modelagem de ADTs, onde os conjuntos de Estados ($Q$) e Eventos ($\Sigma$) de uma FSM canónica serão declarados como Tipos Soma em Haskell. O segundo objetivo será alcançado com a implementação da função de transição pura $\delta$ utilizando \textit{pattern matching} exaustivo, o que é um procedimento metodológico chave para validar a Hipótese de Completude Funcional, demonstrando que o compilador verifica o tratamento de todas as combinações $(\text{Estado} \times \text{Evento})$.

A Fase 2: Gestão do Estado e Sequenciação Monádica endereçará o desafio da pureza na manipulação de estado, cumprindo o terceiro objetivo específico. O estudo da \texttt{State Monad} será traduzido na implementação, garantindo que a sequência de transições seja realizada através de uma função que calcula o novo estado a partir do antigo, mantendo a mudança de estado de forma pura (Milewski, 2020). Esta fase incluirá a sequenciação de eventos utilizando a notação \textit{do-notation} de Haskell, demonstrando como a Categoria Kleisli fornece o formalismo necessário para a composição segura dessas transições monádicas.

A Fase 3: Demonstração da Composabilidade e Segurança é a fase de validação prática da arquitetura. Será selecionado um Estudo de Caso de Domínio que exija a composição de múltiplas FSMs independentes. O procedimento metodológico aqui será compor as máquinas com os construtores do CREM (\texttt{Sequential}, \texttt{Parallel}, \texttt{Alternative} e \texttt{Feedback}), apoiados na composição de Kleisli, e verificar a modularidade arquitetural, simulando o conjunto em malha aberta e em malha fechada e validando o comportamento por enumeração limitada de cenários. O ponto de maior rigor será o Teste de Invalidez de Tipagem, que envolverá a tentativa intencional de codificar uma transição proibida pela lógica do domínio da FSM. A documentação do erro de compilação resultante comprovará a Hipótese de Segurança de Tipos, validando que o sistema de tipos de Haskell é capaz de especificar transições de estado permitidas e proibidas (Perone e Karachalias, 2023) em tempo de compilação.

Por fim, a Fase 4: Análise e Conclusão consistirá na comparação qualitativa dos resultados do modelo funcional puro implementado com o modelo imperativo tradicional, focando na Segurança de Tipos, Completude e Modularidade alcançadas. O código-fonte desenvolvido servirá como a evidência principal para a elaboração final do relatório, que sistematizará o conhecimento adquirido e apresentará a validação conclusiva das hipóteses de pesquisa.

\section{ESTADO DA ARTE}

Esta seção compara quatro maneiras de implementar uma máquina de estados finitos, usando apenas as fontes reunidas até aqui. Não se trata de uma revisão sistemática: as fontes permitem descrever e confrontar as abordagens, mas não afirmar que cobrem a fronteira atual do conhecimento. As lacunas estão listadas no final da seção.

\subsection{O problema nas fontes}

Malakhovski e Shalyto (2009) tratam de autômatos que reagem a eventos, em linguagens funcionais puras. Eles observam que, em linguagens imperativas, o estado costuma ficar em um campo de classe ou em uma variável global, o que não existe em uma linguagem pura. Observam também que os identificadores de estado costumam ser inteiros ou cadeias, às vezes agrupados em um \texttt{enum}, o que não oferece verificação satisfatória quando a função de transição é escrita à mão. Perone e Karachalias (2023) partem de outro incômodo, a distância entre o modelo do domínio e a implementação, comparável à de uma documentação desatualizada, e propõem construir arquiteturas como composição de máquinas de estados que possam ser representadas graficamente.

\subsection{Quatro abordagens}

\textit{FSM imperativa.} O exemplo de Malakhovski e Shalyto (2009, p. 21) é um interruptor de lâmpada em C, com dois \texttt{enum}, uma variável global e uma função com \texttt{switch}. Os autores notam que nada impede, por exemplo, comparar o estado com um evento. A conferência fica com os testes em tempo de execução. Não há, na base, fonte que meça a taxa de erros dessa abordagem em comparação com as demais.

\textit{Haskell ingênuo.} Na versão em Haskell do mesmo exemplo, estados e eventos são tipos algébricos e a função de transição é definida por equações. Os autores afirmam que o compilador confere tipos, completude e ausência de contradições; no exemplo deles, uma equação supérflua gera um aviso. A execução sobre uma lista de eventos é uma dobra à esquerda (\texttt{foldl}). Nos experimentos deste projeto (GHC 9.6.7), a omissão de um par produz um aviso que nomeia o par, e só se torna erro com \texttt{-Werror}. Além disso, uma equação curinga escapa dessa verificação, e o compilador confere se a tabela está completa, não se é a pretendida.

\textit{State Monad.} Os mesmos autores constroem uma mônada própria para o autômato, com operações para ler e gravar o estado, de modo que uma sequência de eventos vira um bloco \texttt{do} (2009, pp. 21--22). Wadler (1992) mostra como mônadas permitem escrever em linguagem pura programas que manipulam estado, com ganho de legibilidade e menos erros do tipo passar o contador errado. Hutton (2016, cap. 12) apresenta a State Monad usada neste projeto. Do lado categorial, Milewski (2020, cap. 4) e Maia (2022) apresentam a composição de Kleisli, Bocchino (2025) relaciona as leis da mônada a essa composição e Jahn (2020) formaliza a mônada como monoide na categoria dos endofuntores. O tipo do estado, porém, não diz quais transições são permitidas: qualquer função de estado em par (valor, novo estado) serve.

\textit{CREM.} Perone e Karachalias (2023) definem a máquina de estados como máquina de Mealy, isto é, uma função com estado que, a cada entrada, emite uma saída e atualiza o estado. As transições permitidas formam uma topologia promovida ao nível dos tipos, e o resultado de cada ação só é aceito se o compilador encontrar a aresta correspondente. As máquinas são folhas de uma árvore, combinadas por construtores de composição. Os autores citam o pacote \texttt{motor}, que também gera diagramas, mas para máquinas isoladas, e apontam limitações: a topologia não se liga às entradas que disparam as transições, e a implementação não explora concorrência.

\subsection{Quadro comparativo}

\begin{table}[h]
\centering
\footnotesize
\begin{tabular}{|p{2.5cm}|p{2.7cm}|p{2.7cm}|p{2.7cm}|p{3.0cm}|}
\hline
 & Imperativa & Haskell ingênuo & State Monad & CREM \\
\hline
Estado & campo ou variável global & argumento e resultado & dentro da computação & dentro da máquina; \texttt{run} devolve a máquina atualizada \\
\hline
Identificador de estado & inteiro, cadeia ou \texttt{enum} & tipo algébrico & tipo algébrico & tipo algébrico promovido a tipo \\
\hline
Equação ausente & testes & aviso do GHC; erro com \texttt{-Werror} & igual ao ingênuo & igual ao ingênuo \\
\hline
Aresta fora do modelo & testes & não detectada & não detectada & rejeitada na compilação \\
\hline
Evento e aresta & código & código & código & só no código \\
\hline
Composição & manual & manual & Kleisli, com leis a provar & construtores do CREM \\
\hline
Diagrama & não previsto & não previsto & não previsto & módulo \texttt{Render} \\
\hline
\end{tabular}
\caption{Comparação das quatro abordagens. As linhas sobre avisos, rejeição de arestas e composição vêm de experimentos deste projeto; as demais, das fontes citadas no texto.}
\end{table}

\subsection{Lacunas}

Faltam fontes, e portanto não há afirmação, sobre: artigos de conferência e periódicos que unam máquinas de estados, Haskell com tipos avançados e teoria das categorias; outras bibliotecas de máquinas de estados em Haskell além das que Perone e Karachalias apenas mencionam; tipos dependentes como alternativa; \textit{statecharts}; verificação de modelos para propriedades de alcançabilidade (os autores do CREM citam TLA+ e Alloy apenas como trabalho futuro); comparações empíricas entre implementações imperativas e funcionais; e requisitos de segurança do domínio de aviônica, que aqui é apenas um domínio de exemplo escolhido pelo projeto. O restante do levantamento será feito junto com o desenvolvimento, conforme o cronograma.

\section{PROGRESSO PARCIAL}

Esta seção apresenta os resultados obtidos até o momento, na ordem cronológica do trabalho, relacionando cada trecho de código ao referencial teórico. Os exemplos iniciais usam uma catraca (\texttt{stateMachine.hs} e \texttt{crem-example}); a partir da seção 5.8, o estudo de caso é a máquina de aviônica de um veículo aéreo não tripulado (UAV), implementada no projeto \texttt{avionics-crem} (repositório Git, versão \texttt{c97981a}). Todo código e toda saída de compilador citados foram produzidos com GHC 9.6.7 e cabal 3.12.1.0, sobre o pacote \texttt{crem} 0.1.0.0.

\subsection{Correspondência entre Tipos de Dados Algébricos e os Conjuntos Q e Sigma}

O primeiro resultado é a tradução direta dos conjuntos finitos de estados ($Q$) e eventos ($\Sigma$) da quíntupla da FSM para tipos soma em Haskell, no arquivo \texttt{stateMachine.hs}:

\begin{verbatim}
data State = Locked | Unlocked deriving (Show, Eq)
data Input = Coin | Push deriving (Show, Eq)
\end{verbatim}

Do ponto de vista categorial, essa codificação instancia um resultado conhecido. Milewski (2020, cap. 5) mostra que o tipo \texttt{Bool}, com dois construtores, corresponde ao número $2$, entendido como o coproduto de duas cópias do objeto terminal ($2 = 1 + 1$). Coproduto é a construção que reúne dois tipos em um só, de modo que cada valor pertence a um ou ao outro; o objeto terminal é o tipo com um único valor. Assim, \texttt{State} e \texttt{Input} são isomorfos a \texttt{Bool}, na categoria Hask, cujos objetos são os tipos de Haskell e cujos morfismos são as funções entre eles. A finitude exigida para $Q$ e $\Sigma$ fica, desse modo, preservada pelo próprio tipo, e não apenas por convenção de nomes.

\subsection{A Função de Transição delta como Morfismo em Hask}

O segundo resultado é a função de transição $\delta: Q \times \Sigma \to Q$, escrita por casamento de padrões:

\begin{verbatim}
transition :: State -> Input -> State
transition Locked   Coin  = Unlocked
transition Locked   Push  = Locked
transition Unlocked Coin  = Unlocked
transition Unlocked Push  = Locked
\end{verbatim}

A assinatura \texttt{State -> Input -> State} é a forma \textit{curried} de uma função cujo domínio é o par \texttt{(State, Input)}. \textit{Currying} é a equivalência entre uma função de dois argumentos, \texttt{(a, b) -> c}, e uma função que recebe um argumento e devolve outra, \texttt{a -> b -> c}; Milewski (2020, cap. 9, seção 9.2) a explica e a liga ao objeto exponencial, e Hutton (2016, seção 3.6) trata o mesmo assunto como funções \textit{curried}. O domínio é o produto categorial de \texttt{State} e \texttt{Input}, construído por Milewski (2020, cap. 5) para pares. A assinatura em Haskell corresponde, portanto, diretamente à definição matemática de $\delta$.

Quanto à completude, o GHC examina as equações e avisa se algum par (estado, evento) ficou sem tratamento. Esse aviso (\texttt{-Wincomplete-patterns}, incluído em \texttt{-Wall}) não impede a compilação por si só; para isso é preciso promovê-lo a erro com \texttt{-Werror=incomplete-patterns}, como faz o projeto \texttt{avionics-crem}. Hutton (2016, seção 3.1) lembra que a inferência de tipos acontece antes da execução, de modo que erros de tipo não ocorrem durante a avaliação; a verificação de exaustividade é uma checagem adicional, também feita em compilação, mas de natureza diferente. Ela confere que cada par tem uma equação, não que a equação seja a correta: trocar \texttt{= Locked} por \texttt{= Unlocked} na segunda linha compila sem aviso.

\subsection{Sequenciação Inicial e Validação Empírica}

A composição de várias transições foi resolvida inicialmente com a função de ordem superior \texttt{foldl}, que aplica $\delta$ repetidamente sobre uma lista de entradas:

\begin{verbatim}
run :: State -> [Input] -> State
run = foldl transition
\end{verbatim}

A execução sobre \texttt{[Coin, Push, Push, Coin, Coin, Push]}, a partir de \texttt{Locked}, produziu o estado final \texttt{Locked}. A função \texttt{transition} é pura no sentido de Milewski (2020): dado o mesmo argumento, devolve sempre o mesmo resultado e não tem efeitos colaterais. Essa pureza, somada à ausência de estado mutável, descrita por Malakhovski e Shalyto (2009), é o que permite aplicar a verificação de tipos e de exaustividade como mecanismos de corretude.

\subsection{Da Composição Ingênua à Necessidade de uma Abstração Monádica}

A composição por \texttt{foldl} basta enquanto o resultado de cada passo é só o novo estado. Quando se quer também uma saída por transição, por exemplo uma mensagem de telemetria, há duas formulações equivalentes.

A primeira não usa mônada. É a máquina de Mealy, em que cada passo tem o tipo $(Q, \Sigma) \to (Q, \text{Saída})$ e a execução sobre uma lista usa \texttt{mapAccumL}, que carrega o estado como acumulador e coleta as saídas. No projeto, essa é a função \texttt{runMealy}.

A segunda é a State Monad, em que o passo é uma função $\Sigma \to ST\;Q\;\text{Saída}$. Nos testes do projeto, as duas formulações foram executadas sobre as mesmas sequências de eventos, a partir dos quatro estados iniciais, e produziram os mesmos pares (saídas, estado final). Logo, produzir saída por transição não exige mônada.

O que a formulação monádica acrescenta é composição. Milewski (2020, cap. 4) chama de funções embelezadas as funções do tipo $A \to M\,B$, em que $M$ é um functor que carrega o efeito (um functor é um construtor de tipos que também sabe aplicar funções a seus conteúdos). A categoria de Kleisli tem os tipos como objetos, essas funções como morfismos e um operador \texttt{>=>}, chamado \textit{fish}, para compô-las; Maia (2022) apresenta o mesmo desenvolvimento a partir do \textit{Writer}. Com ele, um passo posterior pode depender da saída de um anterior sem que o estado seja passado à mão, e o mesmo molde serve para outros efeitos além do estado. A justificativa para a State Monad é, portanto, composição e generalidade, e não necessidade.

\subsection{Reformulação Teórica da FSM sobre a State Monad}

Para tratar do estado, e não apenas de um registro, a FSM foi mapeada sobre a State Monad de Hutton (2016, cap. 12): um \texttt{newtype} que embrulha uma função de estado em par (valor, novo estado). O projeto usa a mesma definição, com o tipo do estado como parâmetro:

\begin{verbatim}
newtype ST s a = ST { runST :: s -> (a, s) }

instance Monad (ST s) where
  st >>= f = ST $ \s ->
    let (x, s') = runST st s
     in runST (f x) s'
\end{verbatim}

As instâncias de \texttt{Functor} e \texttt{Applicative} seguem Hutton e estão em \texttt{Avionics/StateMonad.hs}. Um passo da máquina de aviônica é então:

\begin{verbatim}
avionicsStep :: AvEvent -> ST AvState Output
avionicsStep e = ST $ \q -> (outputOf q e, transition q e)
\end{verbatim}

Esse passo tem o mesmo molde do exemplo de Hutton (2016) para gerar valores frescos, \texttt{fresh = S (\textbackslash n -> (n, n+1))}: um valor interno é passado de forma pura de uma aplicação à seguinte, sem mutabilidade.

\subsection{Sequenciação via Notação do e a Função mapM}

Para generalizar a composição de \texttt{foldl} ao contexto monádico, Hutton (2016, cap. 12) define \texttt{mapM}, que aplica uma função monádica a cada elemento de uma lista e reúne os resultados. No projeto, a função \texttt{runEvents} faz o mesmo para a lista de eventos:

\begin{verbatim}
runEvents :: [AvEvent] -> ST AvState [Output]
runEvents []       = pure []
runEvents (e : es) = do
  o  <- avionicsStep e
  os <- runEvents es
  pure (o : os)
\end{verbatim}

A execução \texttt{runST (runEvents eventos) Grounded} devolve ao mesmo tempo a lista de saídas e o estado final. A notação \texttt{do} é uma forma de escrever expressões com o operador \texttt{>>=} de maneira próxima à das linguagens imperativas, preservando a pureza (Oliveira, 2017).

\subsection{Fundamentação Categorial: o Monad como Monoide na Categoria dos Endofuntores}

Um endofuntor é um functor de uma categoria nela mesma; em Haskell, é um construtor de tipos \texttt{m} com \texttt{fmap}. Uma mônada é um endofuntor $T$ com duas transformações naturais, $\mu : T \circ T \to T$ (achatar) e $\eta : I \to T$ (injetar), cujas leis de associatividade e de unidade têm a forma das leis de um monoide. É a leitura ``monoide na categoria dos endofuntores'', presente em Milewski (2020, cap. 22) e em Jahn (2020, cap. 2).

A consequência prática é que a composição de Kleisli \texttt{>=>} é associativa e tem \texttt{return} como identidade. Isso, porém, só vale se a instância cumprir as leis, e o GHC não as verifica: uma instância que as violasse compilaria normalmente. Por isso as leis foram demonstradas para a instância \texttt{ST}, por raciocínio equacional (Hutton, 2016, seção 16.1; as leis aparecem na seção 12.3), em comentário no arquivo \texttt{Avionics/StateMonad.hs}. Em resumo:

\begin{itemize}
\item identidade à esquerda, \texttt{return a >>= f = f a}: basta substituir as definições e simplificar o \texttt{let} sobre um par conhecido;
\item identidade à direita, \texttt{m >>= return = m}: o \texttt{let} devolve o par \texttt{(fst p, snd p)}, que é igual a \texttt{p};
\item associatividade, \texttt{(m >>= f) >>= g = m >>= (\textbackslash x -> f x >>= g)}: expandem-se os dois lados e achatam-se os \texttt{let} aninhados, obtendo-se expressões idênticas.
\end{itemize}

As demonstrações supõem computações totais e ausência de \texttt{seq}, hipótese análoga à de Hutton (2016, seção 12.4). Bocchino (2025, seções 4.4 a 4.6 e 6.4) mostra que as três leis de Haskell equivalem à associatividade e às identidades da composição de Kleisli, e (seção 7) que, com elas, os tipos e as funções \texttt{a -> M b} formam uma categoria. Além da prova, as leis foram conferidas por exemplo, para os quatro estados e os quatro eventos (sete verificações das leis e duas de consistência, todas aprovadas).

Uma observação sobre o CREM: no código do pacote, a função \texttt{pureResult} exige \texttt{Applicative m}, enquanto \texttt{run} e o sinônimo \texttt{StateMachine} exigem \texttt{Monad m}. \texttt{Applicative} é uma interface mais fraca que \texttt{Monad}, pois oferece \texttt{pure} e \texttt{<*>}, mas não \texttt{>>=}. O CREM é, portanto, parametrizado por um efeito \texttt{m}, o que não implica que a State Monad seja pré-requisito: o CREM nem a usa, e guarda o estado dentro da própria máquina, devolvendo a máquina atualizada a cada \texttt{run}. A relação com a seção anterior é de analogia, não de dependência.

\subsection{Estudo de Caso: Topologia Promovida ao Nível de Tipos}

O framework CREM (Perone e Karachalias, 2023) descreve as transições permitidas por uma topologia: para cada vértice, a lista dos vértices a que ele pode ir. A topologia é um valor comum de Haskell, depois promovido ao nível dos tipos com a biblioteca \texttt{singletons}. Um \textit{singleton} é um tipo que dá, a cada valor de um tipo enumerado, um tipo correspondente no nível dos tipos, de modo que o compilador possa raciocinar sobre ele. A geração desses tipos é feita por Template Haskell, mecanismo de geração de código em compilação.

No estudo de caso de aviônica, os vértices e a topologia são declarados juntos:

\begin{verbatim}
$(singletons [d|
  data AvionicsVertex
    = Grounded | InFlight | EmergencyLanding | Shutdown
    deriving (Show, Eq)

  avionicsTopology :: Topology AvionicsVertex
  avionicsTopology = Topology
    [ (Grounded,         [InFlight, Shutdown])
    , (InFlight,         [EmergencyLanding])
    , (EmergencyLanding, [Grounded])
    ]
  |])
\end{verbatim}

A regra de domínio é que um UAV em voo (\texttt{InFlight}) não pode ser desligado diretamente (\texttt{Shutdown}); por isso não existe essa aresta. Do mesmo modo que as classes de tipos de Hutton (2016), como \texttt{Eq} e \texttt{Ord}, restringem quais tipos servem em um contexto, a topologia restringe quais transições são aceitáveis, tornando a regra de negócio parte do sistema de tipos.

A catraca, usada nas seções anteriores, tem a topologia com as arestas \texttt{Locked} para \texttt{Unlocked} e o inverso. O CREM admite sempre a aresta identidade (um vértice para ele mesmo), por uma instância da classe \texttt{AllowedTransition} (Perone e Karachalias, 2023, Fig. 4); logo, as auto-transições são sempre permitidas. Com dois vértices, as quatro arestas resultantes cobrem todos os pares, e a topologia da catraca é um grafo completo. Por isso ela não serve para demonstrar rejeição.

Também foram escritas, no mesmo projeto, a máquina \texttt{PowerFSM} (subsistema de energia, com os vértices \texttt{Charged}, \texttt{Low} e \texttt{Critical}) e as composições da seção 5.10.

\subsection{Verificação Formal de Segurança de Tipos em Tempo de Compilação}

As ações da máquina de aviônica têm a forma:

\begin{verbatim}
avionicsAction
  :: Applicative m
  => SAvionicsVertex v -> AvEvent
  -> ActionResult m AvionicsTopology SAvionicsVertex v Output
avionicsAction SGrounded TakeoffSignal =
  pureResult MsgTakeoff SInFlight
avionicsAction SInFlight CriticalBattery =
  pureResult MsgEmergency SEmergencyLanding
avionicsAction SInFlight TakeoffSignal =
  pureResult MsgIgnored SInFlight
-- ... (16 pares escritos um a um)
\end{verbatim}

O primeiro argumento de \texttt{pureResult} é a saída; o segundo é o vértice de destino. A função exige a restrição \texttt{AllowedTransition topologia origem destino}, uma classe de tipos que o GHC resolve procurando a aresta na lista da topologia. Eventos que não provocam transição devolvem o mesmo vértice, o que a aresta identidade permite.

Para o teste de invalidez, o módulo \texttt{InvalidInFlightShutdown.hs} tenta a transição proibida, e fica fora do build principal (é um alvo de teste do cabal, construído com \texttt{cabal build invalid-inflight-shutdown --enable-tests}):

\begin{verbatim}
invalidInFlightShutdown SInFlight _ =
  pureResult MsgShutdown SShutdown
\end{verbatim}

A compilação falha. Esta é a mensagem do GHC 9.6.7, sem edição, exceto pelo trecho central omitido, indicado por reticências (o texto integral está no arquivo \texttt{erro\_invalidez\_inflight\_shutdown.txt}):

\begin{verbatim}
invalid\InvalidInFlightShutdown.hs:23:39: error: [GHC-39999]
    * Could not deduce `Crem.Topology.AllowedTransition
                          (Crem.Topology.Topology '[])
                          InFlight Shutdown'
        arising from a use of `pureResult'
      from the context: Applicative m
        bound by the type signature for:
                   invalidInFlightShutdown :: ...
      or from: ()
        bound by a pattern with constructor:
                   SInFlight :: SAvionicsVertex InFlight,
                 in an equation for `invalidInFlightShutdown'
    * In the expression: pureResult MsgShutdown SShutdown
      In an equation for `invalidInFlightShutdown':
          invalidInFlightShutdown SInFlight _
            = pureResult MsgShutdown SShutdown
\end{verbatim}

O trecho \texttt{Topology '[]} mostra como a busca termina: as instâncias de \texttt{AllowedTransition} descartam as arestas da topologia uma a uma e, ao chegar à lista vazia sem encontrar \texttt{InFlight} para \texttt{Shutdown}, o GHC fica sem evidência e rejeita a expressão. A regra de negócio foi, portanto, verificada antes de o programa existir. Isso estende o princípio de segurança de tipos de Hutton (2016) do nível sintático para o nível do domínio, que é a proposta de Perone e Karachalias (2023) de especificar transições permitidas e proibidas no tipo.

Há um limite que precisa constar. A garantia vale para uma transição isolada. O caminho \texttt{InFlight} para \texttt{EmergencyLanding} para \texttt{Grounded} para \texttt{Shutdown} existe na topologia, e é aceito. A regra de que o UAV não se desliga em voo vale para a aresta direta, mas não é propriedade de todas as sequências de eventos; provar propriedades desse tipo exigiria outra técnica.

\subsection{Execução, Validação Cruzada e Composição}

A execução do projeto (\texttt{cabal run avionics-crem}) sobre a sequência de eventos de decolagem, bateria crítica, pouso e corte do sistema produziu a mesma lista de saídas, \texttt{[MsgTakeoff, MsgEmergency, MsgLanded, MsgShutdown]}, e o mesmo estado final, \texttt{Shutdown}, nas três formulações: ingênua, State Monad e CREM. Além disso, as versões ingênua e CREM foram comparadas em todas as 85 sequências de até três eventos a partir de \texttt{Grounded}, e as saídas coincidiram em todas. Essa comparação é uma validação cruzada entre formalizações estruturalmente distintas da mesma máquina.

A composição foi realizada com os construtores do CREM, em \texttt{Avionics/Compose.hs}. Cada folha (\texttt{Basic}) mantém seu próprio estado e sua própria topologia.

\begin{itemize}
\item \texttt{Parallel}: aviônica e energia rodam lado a lado, cada uma com sua entrada e sua saída. Na execução, a entrada \texttt{(TakeoffSignal, Drain)} produziu \texttt{(MsgTakeoff, ReportLow)}.
\item \texttt{Sequential} com \texttt{Alternative}: a saída de energia passa por uma máquina sem estado que a traduz em evento de aviônica, e a aviônica só é acionada quando há evento. Com a aviônica iniciada em \texttt{InFlight} e a bateria em \texttt{Charged}, o segundo \texttt{Drain} levou a bateria ao nível crítico, gerou o evento \texttt{CriticalBattery} e produziu a saída \texttt{MsgEmergency}, isto é, a aviônica passou a \texttt{EmergencyLanding}.
\end{itemize}

O construtor \texttt{Kleisli} não foi exercitado. O \texttt{Feedback} é usado no simulador da seção 5.11.

Os diagramas foram gerados com o módulo \texttt{Render} do CREM, em formato Mermaid (linguagem textual para diagramas): o das topologias de aviônica e de energia e o do fluxo da composição acoplada. Os arquivos estão em \texttt{Obsidian/documentation/} (\texttt{avionics\_topology.mmd}, \texttt{power\_topology.mmd}, \texttt{coupled\_flow.mmd}). Conferiu-se o texto gerado; os diagramas não foram renderizados por uma ferramenta Mermaid.

\subsection{Simulador em Malha Aberta e Fechada e Robustez dos Testes}

Até a seção 5.10, a verificação de comportamento se apoiava em uma sequência fixa de eventos e em 85 sequências curtas. Para tornar o teste mais robusto, foi escrito um simulador em Haskell puro da composição de máquinas (módulos \texttt{Avionics.Sim.*}, sem IO, sem estado mutável e sem aleatoriedade), com uma suíte de testes que percorre todos os cenários até um comprimento dado. Ele valida o comportamento discreto das máquinas e da composição; não modela a física de voo.

Alguns termos. Um \textit{tick} é um passo discreto de simulação, sem correspondência com segundos. Um \textit{traço} é o registro, por tick, da entrada recebida e dos passos de cada máquina (evento, saída e vértice em que ficou). Em malha aberta, cada máquina só recebe o que o cenário manda. Em malha fechada, o ambiente observa a telemetria e realimenta as máquinas, de modo que a saída de uma influencia a entrada da outra.

\textit{Arquitetura.} A planta é a composição em paralelo da aviônica e da energia, com as ações \texttt{avionicsAction} e \texttt{powerAction} originais. O ambiente é uma terceira máquina, \texttt{Mission}, com quatro vértices (\texttt{Standby}, \texttt{CriticalHeld}, \texttt{TouchdownHeld} e \texttt{CutoffHeld}), topologia tipada e 20 equações, sem curinga. Ela é ligada à planta pelo construtor \texttt{Feedback} do CREM: o caminho direto contém a planta e o ambiente, e o caminho de retorno reinjeta os sinais que o ambiente liberou. A State Monad do projeto carrega, de passo em passo, a própria máquina composta como estado. Para relatar em que vértice cada máquina ficou, o simulador observa o vértice de destino dentro da ação, sem reescrever nenhuma regra de transição; o resultado é reconstruído com a mesma prova \texttt{AllowedTransition}, e a garantia por aresta continua valendo.

\textit{Duas decisões de modelagem.} Primeira, o \texttt{Feedback} do CREM é instantâneo: o laço é repetido dentro do mesmo tick até o caminho de retorno parar de emitir. O efeito desejado era influenciar o tick seguinte; ele foi obtido no próprio modelo, com a \texttt{Mission} funcionando como um registrador de um slot: a telemetria relevante a faz reter um evento, e o sinal de relógio (\texttt{Beat}) do tick seguinte a faz liberá-lo. Como o \texttt{Beat} vem só do cenário e a máquina só emite ao recebê-lo, o laço termina; esse argumento está no código, e não há guarda de terminação em execução. Segunda, para que a emergência também afete a energia, foi acrescentado um evento \texttt{ShedLoad} (corte de cargas não essenciais) a \texttt{PowerEvent}. Ele usa apenas a aresta \texttt{Critical} para \texttt{Low}, que já existia, de modo que a topologia de energia não mudou.

\textit{Resultado de compilação.} \texttt{cabal build}, com todos os avisos ligados e a exaustividade promovida a erro, terminou com código 0 e sem avisos do GHC. O teste de invalidez continua falhando com a mesma mensagem (\texttt{Could not deduce AllowedTransition} para \texttt{InFlight} e \texttt{Shutdown}, código de saída 1).

\textit{Traço em malha fechada, cenário de emergência}, obtido com \texttt{cabal run avionics-sim}, na Tabela 2.

\begin{table}[h]
\centering
\scriptsize
\begin{tabular}{|p{0.7cm}|p{3.1cm}|p{4.3cm}|p{3.0cm}|p{2.9cm}|}
\hline
Tick & Entrada ou liberação & Aviônica & Energia & Ambiente \\
\hline
0 & \texttt{TakeoffSignal} & \texttt{MsgTakeoff}, vai a \texttt{InFlight} & \texttt{Charged} & \texttt{Standby} \\
\hline
1 & \texttt{Drain} & \texttt{InFlight} & \texttt{ReportLow}, \texttt{Low} & \texttt{Standby} \\
\hline
2 & \texttt{Drain} & \texttt{InFlight} & \texttt{ReportCritical}, \texttt{Critical} & retém \texttt{CriticalHeld} \\
\hline
3 & libera \texttt{CriticalBattery} & \texttt{MsgEmergency}, vai a \texttt{EmergencyLanding} & \texttt{Critical} & retém \texttt{TouchdownHeld} \\
\hline
4 & libera \texttt{ShedLoad} e \texttt{LandedSignal} & \texttt{MsgLanded}, vai a \texttt{Grounded} & \texttt{ReportLow}, \texttt{Low} & retém \texttt{CutoffHeld} \\
\hline
5 & libera \texttt{SystemCutoff} & \texttt{MsgShutdown}, vai a \texttt{Shutdown} & \texttt{Low} & \texttt{Standby} \\
\hline
\end{tabular}
\caption{Traço do cenário de emergência em malha fechada. As saídas da aviônica são, em ordem, as mesmas quatro da execução em malha aberta.}
\end{table}

O mesmo cenário em malha aberta termina com a bateria em \texttt{Critical} e a aviônica ainda em \texttt{InFlight}: sem realimentação, nada liga uma máquina à outra. Em malha fechada, a bateria termina em \texttt{Low}, porque a emergência dispara o \texttt{ShedLoad}.

\textit{Suíte de testes.} O comando \texttt{cabal test avionics-sim-tests} (comprimento máximo $N = 4$) executou 32 verificações, com 0 falhas. Os cenários de validação cruzada são as sequências de até 4 comandos, isto é, $1 + 4 + 16 + 64 + 256 = 341$; nas verificações com o alfabeto completo de entradas de um tick (20 símbolos), são $1 + 20 + 20^2 + 20^3 + 20^4 = 168\,421$ cenários.

\begin{table}[h]
\centering
\footnotesize
\begin{tabular}{|p{8.3cm}|p{2.6cm}|p{2.2cm}|}
\hline
Verificação & Cenários & Tempo \\
\hline
malha aberta igual à FSM ingênua (saídas e estado final) & 341 & 0,0 s \\
\hline
malha fechada igual à FSM ingênua, com os eventos que a aviônica recebeu de fato & 341 e 168\,421 & 0,01 s e 7,5 s \\
\hline
planta observada igual a \texttt{uavSystem} (Seção 5.10) & 22\,621 & 0,15 s \\
\hline
\texttt{Shutdown} só depois de \texttt{Grounded} (aberta e fechada) & 168\,421 & 0,5 s e 7,5 s \\
\hline
\texttt{EmergencyLanding} só a partir de \texttt{InFlight} (aberta e fechada) & 168\,421 & 0,5 s e 7,4 s \\
\hline
\texttt{ShedLoad} nunca reduz o nível da bateria (aberta e fechada) & 168\,421 & 0,5 s e 7,4 s \\
\hline
determinismo (aberta e fechada) & 168\,421 & 0,7 s e 8,5 s \\
\hline
causalidade: o traço de um prefixo é prefixo do traço (aberta e fechada) & 8\,421 & 0,08 s e 0,9 s \\
\hline
o \texttt{Beat} libera exatamente o que foi retido no tick anterior & 168\,421 & 7,4 s \\
\hline
\end{tabular}
\caption{Verificações da suíte (execução desta revisão; os tempos valem como ordem de grandeza). Em malha fechada compara-se a FSM ingênua com os eventos efetivamente recebidos pela aviônica, porque o ambiente injeta eventos extras.}
\end{table}

\textit{Sanidade das próprias verificações.} Uma propriedade que nunca falha pode ser vacuamente verdadeira. Por isso a suíte inclui um bloco que as aplica a traços corrompidos à mão (por exemplo, um \texttt{Shutdown} vindo de \texttt{InFlight}), e as sete verificações desse bloco confirmam que os defeitos são detectados e que os casos válidos são aceitos. Além disso, nesta revisão foi feito um teste de mutação: trocar, em \texttt{Mission}, o evento liberado por \texttt{CriticalHeld} (\texttt{CriticalBattery} por \texttt{LandedSignal}) fez 6 das 32 verificações falharem; o arquivo foi restaurado em seguida pelo controle de versão. A documentação do repositório registra mutações anteriores, feitas com uma versão anterior da suíte (27 verificações), nas quais 4 e 9 verificações falharam; essas não foram repetidas aqui.

\subsection{Limitações Identificadas e Encaminhamentos}

Registram-se as seguintes limitações.

Primeira, o evento não faz parte do vértice. A topologia restringe apenas quais vértices podem suceder quais, e não qual evento causa cada transição; essa relação fica só no casamento de padrões da ação. Perone e Karachalias (2023, seção 6) registram a mesma limitação como trabalho futuro. Por isso o histórico precisa ser verificado no traço: o tipo não distingue por que se chegou a \texttt{InFlight}.

Segunda, a garantia dos tipos é por aresta (seção 5.9). As propriedades de sequência da seção 5.11 são verificadas por enumeração limitada até $N = 4$, e ausência de contraexemplo nesse recorte não é prova. Além disso, $N = 4$ não alcança a cadeia completa de emergência induzida pela bateria (decolagem, dois \texttt{Drain} e três ticks de liberação exigem seis ticks); esse caminho é coberto apenas pelo cenário fixo da suíte, não pela enumeração.

Terceira, o simulador é discreto: não há física, tempo contínuo nem duração de tick. A semântica do ambiente (bateria crítica vira sensor, emergência vira toque no solo, pouso vira corte) e a do \texttt{ShedLoad} são interpretações do autor do projeto, validadas contra a FSM ingênua e contra as propriedades descritas, não contra requisitos de um domínio real de aviônica; não há fonte na base sobre esses requisitos.

Quarta, o registrador de um slot descarta nova telemetria enquanto está ocupado (quem chega primeiro fica). No recorte $N = 4$ isso não produziu divergência em relação à FSM ingênua, mas a validação não diz se esse comportamento é o desejável; filas ou prioridades são alternativas não exploradas. O ambiente também sempre começa em \texttt{Standby}.

Quinta, o laço não tem guarda de terminação em execução. Depende do argumento descrito na seção 5.11, e uma edição futura de \texttt{Mission} que emitisse em resposta a um sinal realimentado poderia travar o \texttt{Feedback}.

Sexta, o termo ``representável''. No título e no resumo de Perone e Karachalias (2023), uma máquina representável é aquela cuja implementação se mantém sincronizada com sua representação gráfica, ou seja, da qual se gera um diagrama automaticamente. Isso é diferente do functor representável de Milewski (2020, cap. 14): um functor isomorfo ao functor de morfismos $\mathrm{Hom}(a, -)$, isto é, um contêiner indexado por chaves do tipo $a$, de que a tabela de resultados memoizados de uma função é um caso. A ideia de usar o functor representável para indexar a topologia pelo alfabeto $\Sigma$ é uma hipótese de trabalho do projeto, sem sustentação nas fontes da base.

Sétima, o determinismo é trivial em código puro: um traço igual a si mesmo demonstra pouco, e a verificação existe como regressão. A causalidade é o teste mais informativo dessa família.

Oitava, os diagramas Mermaid gerados nesta revisão não foram renderizados; a documentação do repositório informa que versões anteriores dos traços foram renderizadas sem erro com Mermaid 10.9.1, mas o diagrama de emergência foi regenerado depois do \texttt{ShedLoad} e não foi renderizado de novo.

Nona, a hipótese de superioridade sobre as linguagens imperativas só poderá ser avaliada na Fase 4, com a comparação qualitativa prevista; nenhum experimento comparativo foi feito até aqui. O próximo encaminhamento técnico é exercitar o construtor \texttt{Kleisli} e, se for útil, ampliar $N$ ou substituir parte da enumeração por geração de casos com verificação de propriedades.

\subsection{Síntese Geral e Validação Preliminar das Hipóteses}

Na implementação ingênua, a verificação de corretude é a exaustividade do casamento de padrões, que o GHC executa como aviso e que o projeto promove a erro com \texttt{-Werror=incomplete-patterns}. Na etapa monádica, a State Monad e a composição de Kleisli organizam a passagem de estado e saída; as três leis da mônada foram demonstradas para a instância usada e conferidas por exemplo. Na etapa do CREM, a prova \texttt{AllowedTransition} é exigida pelo compilador, e uma transição fora da topologia é rejeitada em compilação, como mostra a mensagem real da seção 5.9. A composição foi realizada com \texttt{Parallel}, \texttt{Sequential}, \texttt{Alternative} e \texttt{Feedback}, e o conjunto foi simulado em malha aberta e fechada, com 32 verificações aprovadas sobre até 168\,421 cenários cada.

Quanto às hipóteses:

\begin{itemize}
\item Completude funcional: confirmada para as funções de transição e de ação da aviônica, da energia e do ambiente, escritas caso a caso e sem equação curinga, com a ressalva de que depende de promover o aviso a erro.
\item Segurança de tipos: confirmada no sentido de que uma aresta ausente da topologia é rejeitada em compilação; não cobre sequências de transições nem a relação entre eventos e arestas. Propriedades de sequência foram verificadas por enumeração limitada, não provadas.
\item Composabilidade e modulação: demonstrada para \texttt{Basic}, \texttt{Sequential}, \texttt{Parallel}, \texttt{Alternative} e \texttt{Feedback}. Falta \texttt{Kleisli}. A propriedade de que a composição preserva as garantias de cada folha foi observada, não formalizada.
\end{itemize}

\section{CRONOGRAMA}

O cronograma a seguir descreve o planejamento de todas as etapas da pesquisa ao longo do período de desenvolvimento do projeto de Iniciação Científica.

\begin{table}[h]
\centering
\small
\begin{tabular}{|p{4.0cm}|p{0.6cm}|p{0.6cm}|p{0.6cm}|p{0.6cm}|p{0.6cm}|p{0.6cm}|p{0.6cm}|p{0.6cm}|p{0.6cm}|p{0.6cm}|p{0.6cm}|}
\hline
\textbf{Etapas} & \textbf{fev} & \textbf{mar} & \textbf{abr} & \textbf{mai} & \textbf{jun} & \textbf{jul} & \textbf{ago} & \textbf{set} & \textbf{out} & \textbf{nov} & \textbf{dez} \\
\hline
Revisão bibliográfica aprofundada & X & X & & & & & & & & & \\
\hline
Modelagem de ADTs e função de transição & & X & X & & & & & & & & \\
\hline
Implementação de State Monad e Kleisli & & & X & X & & & & & & & \\
\hline
Estudo de caso e composição de FSMs & & & & X & X & & & & & & \\
\hline
Testes de invalidez de tipagem & & & & & X & X & & & & & \\
\hline
Análise comparativa e documentação & & & & & & X & X & & & & \\
\hline
Elaboração do relatório final & & & & & & & X & X & X & & \\
\hline
Entrega e apresentação & & & & & & & & & & X & X \\
\hline
\end{tabular}
\end{table}

% ========== REFERÊNCIAS ==========
\newpage
\section{REFERÊNCIAS}

\begingroup
\setlength{\parindent}{0pt}
\setlength{\parskip}{0.6em}
\begin{spacing}{1.0}

ASSOCIAÇÃO BRASILEIRA DE NORMAS TÉCNICAS - ABNT. \textbf{NBR 6023}: Informação e documentação - Referências - Elaboração. Rio de Janeiro, 2018.

ASSOCIAÇÃO BRASILEIRA DE NORMAS TÉCNICAS - ABNT. \textbf{NBR 10520}: Informação e documentação - Citações em documentos - Apresentação. Rio de Janeiro, 2023.

BOCCHINO JR., R. L. \textbf{Monads, Categories, and Computation}. [S.l.]: [s.n.], 2025.

HUTTON, G. \textbf{Programming in Haskell}. 2. ed. Cambridge: Cambridge University Press, 2016.

JAHN, M. \textbf{Monads and their Applications in Haskell}. 2020. Trabalho de Conclusão de Curso (Bachelor's Thesis) -- Department of Mathematics, Technische Universität München, Munich, 2020.

MAIA, D. \textbf{Applied Category Theory: Monads and Haskell}. 2022.

MALAKHOVSKI, Y. M.; SHALYTO, A. A. Konechnye avtomaty v chistykh funktsionalnykh yazykakh programmirovaniya [Máquinas de estados finitos em linguagens de programação funcional pura]. \textit{RSDN Magazine}, São Petersburgo, n. 1, 2009.

MILEWSKI, B. \textbf{Category Theory for Programmers}. [S.l.]: [s.n.], 2020. E-book.

OLIVEIRA, A. G. de. \textbf{Haskell: Uma introdução à programação funcional}. São Paulo: Casa do Código, 2017.

PERONE, M.; KARACHALIAS, G. \textbf{Crème de la Crem: Composable Representable Executable Machines}. 2023. (arXiv:2307.09090).

WADLER, P. Comprehending Monads. \textit{Mathematical Structures in Computer Science}, v. 2, n. 4, p. 461-493, 1992.

\end{spacing}
\endgroup

\end{document}
